% !TEX root = ../main.tex

% 中英文摘要和关键字

\begin{abstract}
  大模型智能体（Large Language Model Agent, LLM Agent）通过工具调用将模型决策作用于文件系统、命令行和外部服务，安全问题随之从回答是否合规转向行动是否被用户真实授权。在实际运行中，风险并不总是来自显式危险请求；不可信环境说明、工具结果误导、模糊指令和长周期任务漂移，都可能使智能体在看似合理的执行过程中偏离用户意图。现有输入过滤、系统提示和单点安全裁判难以稳定刻画行动与授权之间的关系。

  本文研究意图决策不对齐问题，并基于 OpenClaw 与 AgentWard 设计实现运行时防御机制。该机制由用户真实意图分析、意图对齐检测和动态结果处理三部分组成：先抽取任务目标、约束、敏感对象和风险等级，形成结构化授权边界；再在工具执行前比对智能体决策和工具参数是否仍符合该边界；最后根据偏移程度输出放行、审查或阻断。系统以插件形式接入 OpenClaw，并通过时序同步、状态缓存、规则过滤和异常回退保证真实运行中的稳定性。

  本文构建真实 OpenClaw 场景和模拟轨迹两类测试集，在 Qwen3.5-Plus 与 MiniMax-M2.5 上评估方法有效性，并与规则、提示词和单点大模型裁判等基线方法（baseline）比较。真实场景覆盖目标别名重映射和环境诱导删除两类决策偏移攻击及其良性对照；模拟轨迹覆盖对齐、模糊授权、明确不对齐和长周期漂移四类状态。Qwen3.5-Plus 主实验中，原生 OpenClaw 在有效攻击轮次中表现出较高攻击触发率；本文方法在当前测试范围内未观察到攻击残留，并在良性轮次中未造成受保护目标误删或误改。模拟测试结果表明，本文方法对明确不对齐和长周期漂移具有较好的识别能力，但对模糊授权场景仍存在一定不稳定性。实验结果初步表明，意图决策对齐可以作为工具调用前的运行时检查层，为大模型智能体行动安全提供一种可实现的防御思路。

  \thusetup{
    keywords = {大模型智能体, 意图决策对齐, 运行时防御, 工具调用安全, OpenClaw},
  }
\end{abstract}

\begin{abstract*}
  Large Language Model (LLM) agents connect model decisions to file systems, command-line tools, and external services. This shifts safety from whether a response is acceptable to whether an action is authorized. In practice, failures may arise not only from explicit malicious requests, but also from untrusted environmental instructions, misleading tool outputs, ambiguous goals, and long-horizon task drift. Existing input filters, system prompts, and one-shot safety judges have difficulty modeling the authorization relation between a user's intent and an agent's pending action.

  This thesis studies intent-decision misalignment and implements a runtime defense mechanism on OpenClaw and AgentWard. The mechanism first extracts the user's goal, constraints, sensitive objects, and risk level as a structured authorization boundary; then checks before tool execution whether the agent's decision and tool parameters remain aligned with that boundary; and finally maps the result to allow, review, or block. The system is integrated as an OpenClaw plugin with timing synchronization, state caching, rule-based gating, structured outputs, and fallback handling.

  We evaluate the method on real OpenClaw scenarios and simulated trajectories under Qwen3.5-Plus and MiniMax-M2.5, comparing it with rule-based, prompt-policy, and LLM-only baselines. In the Qwen3.5-Plus main experiment, native OpenClaw triggers attacks frequently in valid attack runs, while our method observes no residual attacks within the current test scope and causes no protected-target corruption in benign runs. On 240 simulated trajectories, it performs well on explicit misalignment and long-horizon drift, while remaining less stable on ambiguous authorization cases. These results suggest that intent-decision alignment can serve as a runtime checking layer before tool execution and provide a practical direction for improving LLM agent action safety.

  \thusetup{
    keywords* = {LLM Agent, Intent-Decision Alignment, Runtime Defense, Tool-Call Safety, OpenClaw},
  }
\end{abstract*}
